\chapter{\label{chap:chap3}Proposta}

Este capítulo descreve a solução entregue do \textbf{Dashboard Financeiro}, uma aplicação web voltada ao controle e à análise de finanças pessoais. O sistema já se encontra publicado e combina uma interface moderna, intuitiva e educativa com um backend orientado a serviços capaz de sustentar o crescimento de usuários e integrações.

\section*{Objetivo da Solução}

Oferecer uma ferramenta funcional que atenda a necessidades reais dos usuários e, ao mesmo tempo, estimule o aprendizado e a prática de uma gestão financeira mais consciente. O volume atual demonstra que o sistema completo, educativo e seguro foi efetivamente implementado e encontra-se disponível para uso.

\section*{Organização e Planejamento Metodológico}

A implementação ocorreu de forma incremental e pragmática, priorizando aquilo que agregava mais valor ao produto em cada momento. As entregas eram revisadas diretamente com o orientador e com usuários envolvidos nos testes exploratórios, o que guiou ajustes rápidos sem depender de um board formal ou cerimônias rígidas.

O repositório Git concentrou todas as tarefas, \textit{pull requests} e documentações dos microserviços. Revisões de código, testes locais via Docker Compose e integrações contínuas acompanharam cada incremento, permitindo ajustes rápidos conforme feedback acadêmico e dos usuários envolvidos nos testes exploratórios.

\section*{Funcionalidades Principais}

A aplicação disponibiliza, em produção, as seguintes funcionalidades:

\begin{itemize}
    \item \textbf{Onboarding completo}: cadastro, atualização e exclusão de usuários com validação de CPF/e\hyp mail, autenticação JWT, refresh token, proteção de rotas e preenchimento automático de endereço via integração com a API do ViaCEP no perfil;
    \item \textbf{Gestão financeira diária}: registro de entradas e saídas, categorização flexível (cores e tipos), filtros por período e exportação em tela;
    \item \textbf{Contas bancárias}: múltiplas contas com apelidos, tipos e saldo atualizado automaticamente a cada transação;
    \item \textbf{Metas financeiras}: criação por tipo (\textit{SAVINGS} ou \textit{EXPENSE\_LIMIT}), cálculo de progresso, alocação manual e alertas configuráveis;
    \item \textbf{Dashboards}: visão resumida (cartões, ranking de categorias, metas críticas) e visão analítica (tendência mensal, fluxo diário, distribuição por contas e impacto por categoria);
    \item \textbf{Módulo educacional}: checklists, rituais de aprendizado e dicas contextualizadas com base nos indicadores do usuário;
    \item \textbf{Notificações}: disparo de e\hyp mails via Resend quando metas mudam de status ou novas transações relevantes são registradas;
    \item \textbf{Acesso responsivo}: interface construída com Tailwind e componentes reutilizáveis, adaptada para desktop e dispositivos móveis, sem paywall.
\end{itemize}

Essas funcionalidades compõem o núcleo do produto viável. Evoluções avançadas (integração automática com bancos, análises preditivas e aplicativos móveis) permanecem listadas como trabalhos futuros.

\section*{Definição do MVP}

O MVP referenciado no cronograma corresponde ao conjunto mínimo publicado para validar a proposta com usuários reais. Ele inclui:
\begin{enumerate}
    \item onboarding completo (cadastro, autenticação, verificação e exclusão);
    \item cadastro e gestão de contas bancárias, categorias e transações com atualização automática de saldos;
    \item definição de metas com alertas em tela e via e-mail (metas atingidas/excedidas);
    \item dashboards (resumo + analytics) abastecidos pelos dados reais do usuário;
    \item módulo educacional acoplado às métricas do período.
\end{enumerate}
Esses itens foram priorizados e entregues entre junho e agosto (vide Tabela~\ref{tab:cronograma}); funcionalidades adicionais acrescidas após a validação MVP estão descritas no Capítulo~\ref{chap:chap4}.

\section*{Arquitetura da Aplicação}

A arquitetura organiza-se em serviços especializados que se comunicam por \textbf{APIs REST} e eventos Kafka, seguindo separação de responsabilidades e favorecendo escalabilidade e manutenção.

\subsection*{Componentes Principais}

\begin{itemize}
    \item \textbf{Frontend (React 19 + Vite + Tailwind):} camada de experiência com roteamento protegido, contexto de autenticação e bibliotecas próprias de gráficos;
    \item \textbf{API Gateway (Spring Cloud Gateway):} ponto único de entrada, aplicando autenticação JWT, regras de CORS e roteando chamadas para os demais microserviços;
    \item \textbf{Serviço de Cadastro e Autenticação:} Spring Boot 3.5.5 + Spring Security, responsável por usuários, criptografia de senhas (BCrypt), tokens JWT e publicação de eventos \textit{user.events};
    \item \textbf{Serviço Financeiro}: Spring Boot 3.5.5, Spring Data JPA, Liquibase e Kafka; controla contas, categorias, metas, transações e agregações do dashboard, emitindo eventos \textit{transactions.events} e \textit{goals.events};
    \item \textbf{Serviço de Notificações}: Spring Boot 3.5.5, consumidores Kafka e cliente HTTP para Resend; mantém o cache local de contatos e envia e\hyp mails transacionais;
    \item \textbf{Banco de Dados MySQL}: três esquemas independentes (cadastro, financeiro e notificações) com migrações versionadas por Liquibase.
\end{itemize}

\subsection*{Fluxo de Uso (Exemplo)}

\begin{enumerate}
    \item O usuário acessa a interface (React), que solicita ``\textit{buscar transações do mês}'' para o \textbf{API Gateway};
    \item O gateway valida o \textbf{JWT} (Spring Security) e encaminha a requisição ao \textbf{Serviço Financeiro};
    \item O serviço consulta o \textbf{MySQL} correspondente, recalcula saldos e publica eventos em Kafka quando necessário;
    \item A resposta agregada retorna ao \textbf{Frontend}, que atualiza os gráficos de resumo ou da visão analítica;
    \item Se metas mudarem de status, o evento é consumido pelo \textbf{Serviço de Notificações}, que envia e\hyp mails via Resend com o domínio autenticado.
\end{enumerate}

\subsection*{Diagrama de Arquitetura}

\begin{figure}[H]
\centering
\safeincludegraphics[width=0.95\textwidth]{figs/diagramaarquitetura.png}
\caption{Arquitetura proposta para o \textbf{Dashboard Financeiro}.}
\label{fig:arquitetura}
\end{figure}

\paragraph{Cadastro e Autenticação.}
Criação de conta, autenticação (login), refresh token, gestão de perfil e exclusão são atendidos pelo microserviço \texttt{dashboard-financeiro-cadastro-autenticacao}, que utiliza Spring Security, JWT e MySQL dedicado.

\paragraph{Gestão Financeira.}
CRUD de contas, categorias, transações e metas são tratados pelo microserviço \texttt{dashboard-financeiro}, que centraliza validações, regras de consistência e fechamento mensal, além de publicar eventos Kafka.

\paragraph{Dashboards e Indicadores.}
As consultas agregadas (período, categoria, tendência e metas) são executadas pelo \texttt{DashboardService}, exposto pelo mesmo microserviço financeiro e consumido pelo frontend nas visões ``Quadros'' e ``Analytics''.

\paragraph{Notificações.}
Eventos de usuários, transações e metas são consumidos pelo microserviço \texttt{dashboard-financeiro-notificacoes}, que utiliza Resend para enviar e\hyp mails autenticados com o domínio próprio e persiste logs para auditoria.

\section*{Considerações sobre Distribuição de Funcionalidades entre Componentes}

A separação entre \textit{Serviço de Gestão Financeira} e \textit{Serviço de Dashboard} permite otimizar consultas analíticas sem comprometer a camada transacional. Estratégias como materialized views, índices e endpoints específicos para agregações serão consideradas para garantir desempenho. Autenticação e autorização ficam centralizadas na camada de API para uniformizar políticas de segurança (uso de JWT, roles e scoping mínimo).

O modelo proposto facilita testes unitários e de integração, implantação independente dos serviços e escala horizontal seletiva de componentes com maior demanda (por exemplo, Serviço de Dashboard). Foram implementados testes unitários tanto no frontend (componentes e hooks principais) quanto no backend (serviços de domínio dos microserviços), reforçando a verificação local antes das validações com usuários.
